%%%%%%%%%%%%%%%%%%%%%%%%%%%%%%%%%%%%%%%%%%%%%%%%%%%%%%%%%%%%%%%%%%%%%%%%%%%%%%%%%%
\begin{frame}[fragile]\frametitle{}
\begin{center}
{\Large How ``AI for Teachers'' Programs Are Actually Structured}
\end{center}
\end{frame}

%%%%%%%%%%%%%%%%%%%%%%%%%%%%%%%%%%%%%%%%%%%%%%%%%%%
\begin{frame}[fragile]\frametitle{A Worked Example, Not a Product Pitch}
\begin{itemize}
\item In July 2026, Anthropic launched \textbf{Claude for Teachers}: free premium access for verified US K-12 educators.
\item It is used here purely as a \textbf{worked example} of how a serious education-AI program gets structured, the same four pieces show up under different names across most such programs.
\item Four pieces: a skills library, standards grounding, task automation, and privacy guardrails.
\end{itemize}
\end{frame}

%%%%%%%%%%%%%%%%%%%%%%%%%%%%%%%%%%%%%%%%%%%%%%%%%%%
\begin{frame}[fragile]\frametitle{Piece 1 \& 2: A Skills Library, Grounded in Standards}
\begin{itemize}
\item \textbf{Skills library:} a curated set of ready-made prompts/workflows for recurring teacher tasks, lesson planning, rubric drafting, differentiated explanations, so nobody starts from a blank prompt box.
\item \textbf{Standards grounding:} lesson plans are checked against real curricula (e.g., state academic standards, OpenSciEd, Illustrative Mathematics) rather than the model's generic training knowledge.
\item \textbf{Under the hood, this is RAG}, Retrieval-Augmented Generation: the system looks up the relevant standard first, then writes the lesson plan against it. Same mechanism as the ``search your own documents'' idea from the hands-on section coming up.
\end{itemize}
\end{frame}

%%%%%%%%%%%%%%%%%%%%%%%%%%%%%%%%%%%%%%%%%%%%%%%%%%%
\begin{frame}[fragile]\frametitle{Piece 3 \& 4: Automation and Guardrails}
\begin{itemize}
\item \textbf{Task automation:} agentic tools that run a recurring job on a schedule, e.g. a daily 4pm exit-ticket review, without a teacher re-prompting every day.
\item \textbf{Privacy guardrails:} FERPA-aligned terms, co-developed with a practitioner body (in this case, the American Federation of Teachers), not written by engineers alone.
\item \textbf{The pattern to look for} in any tool your school considers: does it know your curriculum, can it run on autopilot for the boring parts, and does it say clearly what it does with student data?
\end{itemize}
\end{frame}

%%%%%%%%%%%%%%%%%%%%%%%%%%%%%%%%%%%%%%%%%%%%%%%%%%%
\begin{frame}[fragile]\frametitle{This Is Not a Niche Effort}
\begin{itemize}
\item Beyond the K-12 program, the same company partnered with \textbf{Teach For All} to bring AI training to educators across \textbf{63 countries}.
\item Over \textbf{100,000 teachers and alumni} in that network alone have access to develop AI fluency and adapt these tools to their own classrooms.
\item The scale is the signal: this is not a pilot program anymore, it is infrastructure being built under the profession.
\end{itemize}
\end{frame}
